% CLAIM: stated, not buried. docs/05's list plus the earned ones.
\section{Limitations}
\label{sec:limitations}

\begin{itemize}
  \item \textbf{Eight devices, one host, a fat fabric.} The GPU results are
    a single NVLink node and the TPU results a single v5e host: the two
    interconnects on which a model-sized all-reduce is cheapest. The scalar
    folk claim was made for commodity clusters of hundreds of workers
    \citep{salimans2017}, where it may well hold; the dense-side result here
    is the easy case for B. Scaling past a host boundary is untested and may
    move the crossover substantially.
  \item \textbf{Costs are attributed in bytes, not in time.} Every timing
    is end-to-end wall clock; no collective is timed in isolation at $P$
    bytes on either fabric, no achieved bandwidth is reported, and the
    regeneration-versus-matmul-versus-collective split of a generation is
    not profiled. The crossover is an observation, not a cost model.
  \item \textbf{The placement space is larger than A and B.} Reduce-scatter
    with sharded update application, compressed collectives, hierarchical
    contraction, and overlapping the all-reduce with the next evaluation are
    all unmeasured (Section~\ref{sec:contraction}); overlap in particular
    could remove the low-rank side of the crossover.
  \item \textbf{The systems workload is a simplified block.} The scaling and
    cost grids use square matrices, $d_{\mathrm{ff}}{=}d$, one head, no
    learned norms, MSE loss, batch 8, and sequence length 32. These choices
    favor attribution over realism, so absolute ratios should be read as
    properties of this workload. The crossover recurs on a real model with
    both signs intact (Table~\ref{tab:tb5}), but that is a narrower anchor than
    the grid.
  \item \textbf{Models must adopt the library's seams.} Parameterized matmuls
    and embedding lookups must route through \texttt{shardes.nn}; an existing
    model does not run unmodified, and unsupported layer types require a new
    seam. This is the API cost of never materializing low-rank noise.
  \item \textbf{The crossover is bracketed at sweep resolution.} The $(N,d)$
    grid behind Figure~\ref{fig:f2} is $3\times2$ per platform, with holes at
    the memory ceiling. It establishes regions and sign changes, not a precise
    contour.
  \item \textbf{Absolute numbers are platform-specific.} The v5e has 16\,GB
    per chip versus the A100's 80\,GB, and no throughput number here should be
    transferred to another accelerator or interconnect without measurement.
    The real-model crossover was measured only on the v5e; cross-accelerator
    comparisons use the controlled block workload.
  \item \textbf{Task validation is one task and one training model size,
    with a weak positive control.} Countdown with a 0.5B-parameter model,
    three seeds, and $N{=}30$ cannot
    establish rank parity on other tasks, populations, or scales; the
    setup saturates early and resolves a frozen arm but was not shown to
    resolve a modestly worse one. The
    alignment chain spans 0.5B and 1.5B models, but no 1.5B training run was
    performed.
  \item \textbf{The low-rank correction is empirical.} The 0.53 factor has a
    cell-to-cell scatter of about a quarter and drifts to 0.40 at 1.5B; it
    patches a sampling-dimension variable that counts draws rather than
    measuring overlap with the gradient, and no derivation of the right
    variable is offered.
  \item \textbf{NLL alignment is a proxy, not the task objective.} The exact
    gradients in Tables~\ref{tab:tb4} and~\ref{tab:tb6} are for
    teacher-forced NLL, whereas the campaign optimizes a nondifferentiable
    reward from decoded answers. The task outcome is consistent with the
    alignment prediction but cannot establish the same mechanism.
\end{itemize}
